\emph{Step 1 (the index).} $\ell\mid k_n=[\RE:A_n]$ (\cite{KY} Eq.~(17), an exact equality).

\emph{Step 2 (types).} Write $[r_a]:=r_aA_n$ for the coset of the real formal root $r_a\in A_n^{1/\ell}$ (a number) in the group $A_n^{1/\ell}/A_n$; the identification \eqref{eq:ident} sends $\bar a\in R_n/\ell R_n$ to $[r_a]$, and it is well defined because replacing $a$ by $a+\ell b$ multiplies $r_a$ by $u_b\in A_n$ (the real $\ell$-th root is unique since $\ell$ is odd: \path{WeberRoots.odd_pow_injective}, \path{WeberRoots.root_change_rep}, \path{WeberRoots.coset_eq_of_change_rep}). Since $A_n\subseteq\RE$ (\cite{KY} Def.~2.1), the set $\{[\varepsilon]:\varepsilon\in\RE\cap A_n^{1/\ell}\}$ is a subgroup of $A_n^{1/\ell}/A_n$; it is what \cite{KY} write $\RE/A_n$.

\emph{Step 3 (the published input).} \cite{KY} Prop.~4.1 supplies an irreducible $f_0$ with $\{g\cdot\varepsilon_n: g\in\Z[x],\ g\bmod(x^m+1)\in M_{f_0}\}\subseteq\RE/A_n$, where $g\cdot\varepsilon_n$ is the real root $r_g$ and the set is a set of cosets; in the identification \eqref{eq:ident} this says that every coset $[r_a]$ with $\bar a\in M_{f_0}$, i.e.\ with $a\in L_{f_0}$, belongs to the subgroup of Step~2: $[r_a]=[\varepsilon]$ for some $\varepsilon\in\RE\cap A_n^{1/\ell}$. (We do not say that a coset is a subset of $\RE/A_n$; membership of the coset in the image of $\RE\cap A_n^{1/\ell}$ is the statement.)

\emph{Step 4 (from the coset to the number).} $[r_a]=[\varepsilon]$ means $r_a=\varepsilon c$ with $c\in A_n$. Since $A_n\subseteq\RE$ and $\RE$ is a group, $r_a\in\RE$. This holds for every $a\in L_{f_0}$; if $a$ is replaced by $a+\ell b$, $r_a$ is multiplied by $u_b\in A_n\subseteq\RE$, so membership in $\RE$ is preserved, as it must be since $\ell R_n\subseteq L_{f_0}$.

\emph{Step 5 (the log vector).} $r_a$ is the real $\ell$-th root of $u_a$, so $\ell\log|\sigma^j r_a|=\log|\sigma^j u_a|=(H_na)_j$ for every real embedding $\sigma^j$, $j<2m$: the log vector of $r_a$ is $\ell^{-1}H_na$. The base branch $a\in\ell R_n$ is Lemma~E$'$.
